% ID: SEL-D0046-OA-E04
% ESTADO: cotejada. Segunda pasada visual y borrador compilado revisados.
% Fuente: Selectividad/sel_2014_matematicas_aplicadas/Matematicas_CCSS_JUN.pdf
% SHA-256: df7375f8183eafb6f45e9c65083e68d528ffabe07d66717b0ccf0f020c5b064f
% Página: 1; opción A; ejercicio original 4.

Se quiere hacer un estudio de mercado para conocer el precio medio de los libros de narrativa que se venden en la actualidad. Para ello se elige una muestra aleatoria de 121 libros, encontrando que tienen un precio medio de 23 \texteuro{}. Se sabe que el precio de los libros de narrativa sigue una distribución Normal con media desconocida y desviación típica 5 \texteuro{}.
\begin{enumerate}[label=\alph*)]
\item \textbf{(1.5 puntos)} Obtenga un intervalo de confianza, al 98.8\%, para el precio medio de esos libros.
\item \textbf{(1 punto)} \textquestiondown{}Cuántos libros habría que elegir como muestra para que, con la misma confianza, el error máximo de la estimación no excediera de 1\texteuro{}?
\end{enumerate}
